\documentclass{book}
\usepackage{hyperref}
\usepackage[margin=1in]{geometry}
\usepackage{ulem}
\usepackage{tipa} % Add this line to support \textipa command
\hypersetup{
    colorlinks=true,
    linkcolor=black,
    filecolor=magenta,      
    urlcolor=cyan,
    }

\urlstyle{same}

\title{French Conjugation}
\author{Junzhang Luo}

\begin{document}

\tableofcontents

\chapter{DELF A1}

    \section{Sample Paper 2}

        \begin{enumerate}
            \item Vous venez de vous installer dans une nouvelle ville pour vos études/votre travail. Vous écrivez un e-mail à un(e) ami(e) pour lui raconter votre nouvelle vie.
                    Vous lui expliquez ce que vous faites et vous lui demandez de ses nouvelles. \\
                    Objet: Nouvelle vie à Montréal \\
                    Bonjour mon ami, 
                    J'ésprait que vous allez bien. J'ai arrivé à Montréal en train hier et je suis à le marché des fenetres. J'ai en train de chercher les fênetres pour mon nouvelle apartement. 
                    C'est suite à côte de centre-ville, environ deux kilomètre. Qu'est-ce que tu vas recénment? Je voudrais connatire. \\
                    Merci,\\
                    Alex 

                    \medskip
                    \textbf{Annotations (Version Corrigée):}
                    \begin{itemize}
                        \item \sout{J'ésprait que vous allez bien.} J'espère que tu vas bien.
                        \item \sout{J'ai arrivé à Montréal en train hier} Je suis arrivé à Montréal en train hier soir
                        \item \sout{je suis à le marché des fenetres. J'ai en train de chercher les fênetres pour mon nouvelle apartement.} je m'installe peu à peu. En ce moment, je cherche des fenêtres pour mon nouvel appartement et je compare les prix dans plusieurs magasins du quartier.
                        \item \sout{C'est suite à côte de centre-ville, environ deux kilomètre.} L'appartement est tout près du centre-ville, à environ deux kilomètres, ce qui est pratique pour les courses et le métro.
                        \item J'aime beaucoup l'ambiance ici: les rues sont animées, et il y a déjà des marchés d'automne avec des produits locaux. Ce week-end, je compte visiter le Mont Royal et essayer un café dont on m'a parlé.
                        \item \sout{Qu'est-ce que tu vas recénment? Je voudrais connatire.} Et toi, comment vas-tu? Donne-moi de tes nouvelles quand tu as un moment. Si tu passes par Montréal, on pourrait se voir et découvrir un resto ensemble!
                        \item \sout{Merci,} Cordialement,
                    \end{itemize}
                    Alex
        \end{enumerate}
    \newpage
    \section{Sample Paper 3}

        \begin{enumerate}
            \item Vous répondez à votre ami: vous acceptez l'invatation, vous annoncez votre arrivée, vous proposez quelque chose (40 à 50 mots) \\
                    Bonjour Montpellier,
                    Merci pour ton invitation, j'éspiere que tu vas bien. J'accepte l'invitation et j'étais à la caompagne dans le samedi 25 avec ma femme. Je vais apporter une gateau, des biscuits et les pains du raisin. Aussi, je pense que on devrait avoir un feu de camp pour griller des guimauves.
                    J'ai hâte de te voir. \\
                    Cordialement, \\
                    Alex \\
                    \medskip 
                    \textbf{Annotations (Version Corrigée):}
                    \begin{itemize}
                        \item \sout{Merci pour ton invitation, j'éspiere que tu vas bien. J'accepte l'invitation} Merci pour ton invitation - je l'accepte avec grand plaisir.
                        \item \sout{j'étais à la caompagne dans le samedi 25 avec ma femme.} J'arriverai le samedi 25, vers 16 h, probablement en conduisant.
                        \item \sout{Je vais apporter une gateau, des biscuits et les pains du raisin.} Je peux apporter un gâteau, des biscuits et des pains aux raisins.
                        \item \sout{Aussi, je pense que on devrait avoir un feu de camp pour griller des guimauves.} On pourrait aussi faire un feu de camp pour faire griller des guimauves ou autres choses, si tu es d'accord.
                        \item J'ai hâte de te voir.
                        \item \sout{Cordialement,} À très bientôt,
                    \end{itemize}
                    Alex
        \end{enumerate}
    \newpage
    \section{Sample Paper 4}
        \begin{enumerate}
            \item Vous écrivez une lettre à un(e) ami(e) français(e) pour l'inviter dans votre pays pendant les vacances. Vous lui parlez des activités que vous pouvez faire ensemble. \\
                    Bonjour mon ami, \\
                    J'éspiere que tu vas bien. J'invite tu visites mon pays pendant ton vacance si tu as les temps. On peut visiter au parc, ses promenade dans la rue pour faire du shopping, et finalement, 
                    je peux essayer les répas traditionale, comme les desserts, le jus de la prune aigre et plus. Aussi, il y a une fête du culturel jusqu'à l'automne. Il y aura plusieurs spectacles interprétés par des locaux. 
                    Je pense que serais amusante et spectaculaire. \\
                    J'ai hâte de te voir. \\
                    Cordialement, \\
                    Alex
                    \medskip 
                    \textbf{Annotations (Version Corrigée):}
                    J'espère que tu vas bien. Je voulais t'inviter à venir visiter mon pays pendant les prochaines vacances, si tu as le temps. Ce serait une belle occasion de passer du temps ensemble et de te faire découvrir ma région. \\
                    Nous pourrons visiter le grand parc national, nous promener dans les rues animées pour faire un peu de shopping. Ensuite, nous pourrons goûter aux repas traditionnels: des plats typiques, des desserts locaux et même le jus de prunes aigres, une spécialité de chez nous. \\
                    De plus, il y aura une fête culturelle qui continuera jusqu'à l'automne tous les samedi. Tu pourras assister à plusieurs spectacles interprétés par des artistes locaux: musique, danse et théâtre. Je suis sûr que ce sera amusant et spectaculaire. \\
                    J'ai vraiment hâte de te revoir et de partager ces expériences avec toi. \\
                    Amicalement, \\
                    Alex
        \end{enumerate}

\chapter{DELF A2}

    \section{Sample 1}
        \subsection{Exercice 1}
            J'ai atterri à l'aéroport de Paris dans le soir. J'étais fatigué mais trés excitemment, hâtais des prochaines jours dans Paris.
            Le matin prochain, j'ai visité le Château de Versailles, que est un palais renommé mondiale. 
            J'ai passé plusieurs heures là, apprécié les sculpteurs, l'œuvres d'art différents.
            J'imprimais par le technique d'art et le structure d'immeuble. 
            Après avoir essayé les repas locales, je me suis promenadé le long de la rive du fleuve de La Seine.
            J'ai visité la Cathédrale de Paris dans le derniere jour avant partir pour retourne. 
            Géneralement, le vie citadine en Paris donne un ambiance détendu et relaxé. J'ai apprécié mon séjour dans Paris.
        
        \subsection*{Fixed}
            Jour 1 \\
            Je suis arrivé à l'aéroport de Paris \textcolor{red}{le soir}. J'étais fatigué, mais très enthousiaste. J'avais hâte de découvrir la ville dans les jours suivants. \\
            Jour 2 \\
            Le lendemain matin, j'ai visité le château de Versailles, \textcolor{red}{un palais mondialement connu}. J'y ai passé plusieurs heures et j'ai admiré les sculptures ainsi que les différentes œuvres d'art.
            J'ai été impressionné par la richesse des détails, la technique artistique et l'architecture du bâtiment.
            \\ Jour 3 \\
            Après avoir goûté des plats locaux, je me suis promené le long des quais de la Seine. 
            L'ambiance m'a plu: malgré la foule, j'ai trouvé les Parisiens plutôt calmes, et la ville semblait pleine de vie.
            \\ Jour 4 \\
            Avant de repartir, j'ai visité la cathédrale Notre-Dame (de l'extérieur). 
            J'ai ressenti une certaine émotion en pensant à son histoire. 
            Globalement, la vie à Paris m'a donné une impression dynamique mais agréable. J'ai beaucoup apprécié mon séjour et j'aimerais y retourner.

        \subsection*{Remarks}
            \begin{itemize}
                \item plaire à quelqu'un instead of aimer. Ça me plaît (I like it).
                \item ignore ``in'' for time like evening.
            \end{itemize}
        
        \newpage
        \subsection{Exercice 2}
            Salut, \\
            Merci beaucoup pour l'invitation et j'espèrais que tu va bien. 
            J'aimerais aussi bien te voir, cependant, je n'ai pas assez de temps pour aller à Paris parce que je vais au Étas-Unis avec mes amis.
            Si tu n'dérange pas, nous pouvons nous rencontrer avant les vacances et avoir une conversation ensemble.

        \subsection*{Fixed}
            Salut, \\
            Merci beaucoup pour ton invitation! J'espère que tu vas bien. Moi aussi, j'aimerais te voir.
            Malheureusement, je ne pourrai pas venir à Paris pendant les vacanes, car je pars aux États-Unis avec des amis. 
            En revanche, on pourrait se rencontrer avant les vacances, par exemple pour prendre un café et discuter. Sinon, je peux t'appeler vers 20 h 30 cette semaine.
    \newpage
    \section{Sample 2}
        
        \subsection{Exercice 1}
            J'ai passé une nuit blanche extraordinaire dans le 2 octobre derneir. Le musique était très bonne et j'ai apprécié l'atomosphère vivante.
            Les gens que j'ai parlé est très gentil. J'ai vu les auteurs jouant des spectacles de culture différent. 
            Il était très intérésant à apprendre et être exposer. L'événement était choquant et une toute nouvelle expérience pour moi.
            Je n'ai pas aimé le repas à l'événement parce qu'il était principalement de la restauration rapide pour vente.
            Globalement, il était une expérience intéréssant cependant, je suis trés fatigué et je pense que je n'serai pas les autures événement où je dois rester éveillé.

        \subsection*{Corrected Version}
            Le 2 octobre dernier, j'ai passé une nuit blanche extraordinaire dans la rue. La musique était très bonne et j'ai beaucoup apprécié l'atmosphère animée.
            Les gens à qui j'ai parlé étaient très gentil et chaleureux. J'ai vu des artistes présenter des spectacles de cultures différents. 
            C'était très intéresent, car j'ai pu découvrir de nouvelles formes d'art et être exposé à plusieurs traditions culturelles. \\
            Cet événement a été une expérience surprenante et totalement nouvelle pour moi. 
            En revanche, je n'ai pas aimé la nourriture, parce qu'il y avait surtout de la restauration rapide à vendre. \\
            Globalement, c'était une expérience très intéressante. Cependant, j'étais très fatigué à la fin, et je ne pense pas participer à d'autres événements où il faut rester éveillé toute la nuit.
        
        \subsection*{Remarks}
            \begin{itemize}
                \item No ``dans'' for date.
                \item use ``qui'' for parler à quelqu'un.
                \item présenter un spectacle, offrir un spectacle, jouer dans un spectacle, etc. instead of ``jouer des spectacles''.
            \end{itemize}

        \subsection*{B2 Level Correction}
            Le 2 octobre dernier, j'ai passé une nuit blanche extraordinaire dans les rues de la ville. 
            Dès le début de la soirée, j'ai été séduit par l'ambiance festive et animée. La musique était excellente et on pouvait entendre des styles très variés selon les endroits.
            \\
            J'ai également beaucoup apprécié les spectacles de rue. J'ai vu des artistes présenter des performances inspirées de cultures différentes, ce qui m'a permis de découvrir de nouvelles formes d'expression artistique.
            Les gens à qui j'ai parlé étaient très gentils, ouverts et chaleureux, ce qui a rendu l'expérience encore plus agréable.
            \\
            En revanche, je n'ai pas vraiment aimé la nourriture proposée, car il y avait surtout de la restauration rapide.
            J'aurais préféré des plats plus originaux ou plus typiques.
            \\
            Globalement, cette nuit blanche a été une expérience inoubliable, à la fois surprenante et enrichissante. 
            Même si j'étais très fatigué à la fin, je garde un excellent souvenir de cette soirée grâce à l'ambiance, aux spectacles et à l'énergie des gens.
        \newpage
        \subsection{Exercice 2}
            Bonjour Pierre, pour le destination du voyage dans l'été, je vais recommander le parc nationnel de Banff. Il est situé le cœur de montagne Rockie.
            Le parc est massif, il y a plusieur routes qui tu peux choissir parmi. De plus, il y a diverses espèces d'animaux, tu peux voir l'ours, l'élan par exemple sur la côté de la rue.
            Le plus important, les lacs sont spectaculaires et chaque d'eux a traits distinctifs, certains sont bleus en revanche certains sont verts. 
            C'est un place idéal pour apprécie de la nature, un place idéal pour vacances d'été.

        \subsection*{B2 Level}
            Bonjour Pierre, 
            \\
            Pour tes vacances d'été, je te recommande le parc national de Banff, au Canada. 
            Il est situé au cœur des montagnes Rocheuses et c'est un endroit magnifique pour les amoureux de la nature.
            \\
            Le parc est immense et il y a plusieurs routes et sentiers que tu peux choisir selon tes envies. 
            Pendant mon voyage, j'ai fait de la randonnée, j'ai admiré les lacs et j'ai pris beaucoup de photos. 
            On peut aussi y voir différentes espèces d'animaux, par exemple des ours ou des élans au bord de la route.
            \\
            Ce que j'ai surtout aimé, ce sont les lacs, car ils sont spectaculaires. Chacun a ses propres caractéristiques: certains sont d'un bleu éclatant, tandis que d'autres sont plutôt verts.
            C'est donc un endroit idéal pour profiter de la nature et passer des vacances d'été inoubliables.
    \newpage
    \section{Sample 3}
        \subsection{Exercice 1}
            Le Samedi dernier, j'ai mangé diner avec la grande famile du côté de ma mère. Il y avait environ douze gens. 
            C'était mon premiere fois recontrer les gens majoritaires là. 
            J'étais trop surpris que les familes du ma mère ont relatifs nombreux.
            Tout le monde ont arrière-plan différent, comme les professeur, manageur de l'usine de le papier.
            J'ai particulèrement apprécié la ambiance de retrouvailles. Le repas offre était très delicieux.
            
        \subsection*{B2 Correction}
            Samedi dernier, j'ai dîné avec la grande famille du côté de ma mère. Il y avait environ douze personnes.
            \\
            C'était la première fois que je rencontrais la majorité d'entre elles, donc j'étais très surpris de voir que la famille de ma mère était si nombreuse. \\
            Tout le monde avait un parcours différent: il y avait, par exemple, des professeurs, un directeur d'usine de papier et présentateur de télévision.
            \\
            J'ai particulièrement apprécié l'ambiance chaleureuse de ces retrouvailles. De plus, le repas qui nous a été servi était délicieux.
            \\
            Ce moment a été inoubliable pour moi, car j'ai eu l'impression de mieux connaître ma famille et de me rapprocher d'elle.
        
        \subsection*{Notes}
            \begin{itemize}
                \item personnes instead of gens (used for general)
                \item la première fois que + verb
                \begin{itemize}
                    \item fois is feminine
                \end{itemize}
                \item être surpris de voir que
                \item use ``de'' for possession
                \item ``si'' means ``so'' in ``était si nombreuse''
            \end{itemize}

        \newpage 

        \subsection{Exercice 2}
            Salut Mathilde,
            \\
            Remercie beaucoup pour l'invitation et je l'accepte. J'espère qu tu vas bien aussi. 
            J'a acheté le billet et le vol atterrira à Paris dans le 20 juillet 2026 environ à dix-huit heures moins le quart.
            Je pense que le douane prendra une demi-heure et je serai hors de l'aéroport par dix-neuf heures. 
            Je te recontrai à ce moment-là, appelle moi ou envoyez moi un message quand t'arrivé. 
            J'ai hâte de te voir et passer le temps avec toi.
            \\
            Sincèrement,
            Ann

        \subsection*{Corrected Version}
            Salut Mathilde,
            \\
            Merci beaucoup pour ton invitation, que j'accepte avec plaisir. J'espère que tu vas bien aussi.
            \\
            J'ai déjà acheté mon billet d'avion. Mon vol arrivera à Paris le 20 juillet 2026 à 17 h 45. 
            Je pense que je mettrai environ une demi-heure à passer la douane, donc je serai à la sortie de l'aéroport vers 18 h 15.
            \\
            Pendant mon séjour en France, j'aimerais beaucoup visiter quelques monuments célèbres, goûter la cuisine française et me promener dans la ville avec toi.
            J'amerais aussi découvrir la fête nationale et participer aux célébrations.
            \\
            J'ai hâte de te voir et de passer du temps avec toi.

        \subsection*{Notes}
            \begin{itemize}
                \item Don't use ``dans le'', use ``le'' for date.
                \item Retrouver is better than recontrer, since not first time.
            \end{itemize}

        \subsection*{C1 Version}
            Merci infiniment pour ton invitation, que j'accepte avec beaucoup de joie. 
            Cela me fait très plaisir de savoir que je pourrai passer quelques jours chez toi à l'occasion de la fête nationale.
            J'espère que tout va bien pour toi et que les préparatifs se passent comme tu le souhaites.
            \\
            J'ai déjà réservé mon billet, et mon avion atterrira à Paris le 20 juillet 2026 à 17 h 45. 
            Si tout se passe normalement, je devrais avoir passé la douane et récupéré mes bagages vers 18 h 15. 
            Nous pourrons donc nous retrouver à la sortie de l'aéroport à ce moment-là.
            \\
            Pendant mon séjour, j'aimerais non seulement assister aux célébrations de la fête nationale et admirer les feux d'artifice, 
            mais aussi découvrir davantage la culture française. J'aimerais visiter quelques musées ou monuments emblématiques, flâner dans différents quartiers, 
            profiter de l'ambiance des cafés et, bien sûr, goûter à plusieurs spécialités locales. Ce serait également l'occasion idéale de passer du temps ensemble et de mieux connaître ta ville à travers tes yeux.
            \\
            Je me réjouis déjà de ce voyage et de nos retrouvailles. Merci encore pour ton accueil si généreux.

    \newpage
    \section{Sample 4} 
        \subsection{Exercice 1}
            Le grand festival que la famille célébere seulment est la fête du printemps, 
            òu nous inviteraions autre membre du famille et les amis à venir chez moi pour dinner.
            Généralement, on leur parler et en faisant le diner, qui est grand et copieux. 
            Le repas typique est les boulettes ou les boules des riz selon la région que vous venez.
            Pour ma famille, on vient du nord de la Chine, donc c'est les boulettes pour nous.
            C'est un effort collaboratif pour faire les boulettes, òu tout le monde parle en les faisant.
            Après dinné, les activités inclus mahjong, jouer aux cartes etc. 
            \\
            Je pense que cet forum est une idée excellente pour apprendre les fêtes dans famille différents.

        \subsection*{Correction}
            La grande fête que ma famille célèbre surtout est la fête du Printemps. 
            \\
            À cette occasion, nous invitons d'autres membres de la famille ainsi que des amis à venir chez nous pour dîner.
            \\
            En général, on parle beaucoup ensemble tout en préparant le repas, qui est souvent copieux et convival.
            \\
            Le plat typique est composé de boulettes ou de boules de riz, selon la région d'où l'on vient.
            \\
            Dans ma famille, nous venons du nord de la Chine, donc, pour nous, ce sont les boulettes qui occupent une place centrale.
            \\
            La préparation des boulettes est un véritable effort collectif, pendant lequel tout le monde discute et partage un bon moment.
            \\
            Après le dîner, nous jouons souvent au mah-jong ou aux cartes.
            \\
            Personnellement, je pense que ce forum est une excellente idée, car il permet de découvrir les traditions familiales de différentes cultures.
            
        \subsection*{C1}
            Dans ma famille, la célèbration la plus manquante est sans aucun doute la fête du Printemps,
            qui occupe une place centrale dans notre vie familiale. 
            À cette occasion, nous réunissons non seulement les proches de la famille, mais aussi parfois des amis très chers,
            afin de partager un grand repas dans une atmosphère chaleureuse et conviviale.
            \\
            Ce que j'apprécie particulièrement, c'est le caractère collectif de la préparation.
            Dans ma famille, qui est originaire du nord de la Chine, les boulettes constituent le plat traditionnel incontournable. 
            Leur confection devient presque un rituel: chacun met la main à la pâte, tout en discutant, en échangeant des souvenirs ou simplement en profitant du moment passé ensemble.
            Une fois le repas terminé, la soirée se prolonge généralement autour de jeux comme le mah-jong ou les cartes,
            ce qui renforce encore davantage le sentiment de proximité entre les membres de la famille.
            \\
            À mes yeux, les fêtes de famille jouent un rôle essentiel, car elles permettent de transmettre des traditions,
            de resserrer les liens affectifs et de créer des souvenirs durables. C'est pourquoi je trouve ce thème particuilèrement intéressant: il révèle à quel point les coutumes familiales,
            bien qu'elles diffèrent d'une culture à l'autre,
            répondent souvent au même besoin de partage et d'appartenance.

        \newpage
        \subsection{Exercice 2}
            Merci infiniment pour ton invitation, félicitations pour votre mariage. 
            J'ai accepté ton invitation avec beaucoup de plaisir. 
            Je arriverais avec mon ami et j'ai hâte de témoigner ce moment le plus important pour vos.
            N'hésitez pas à me faire savoir si vous avez rien que je peux vous aider avec installer le lieu.

        \subsection*{Correction}
            Merci infiniment pour votre invitation, et toutes mes félicitations pour votre mariage.
            J'accepte voitre invitation avec grand plaisir. 
            \\
            Je viendrai avec mon ami et j'ai hâte d'être présent pour ce moment si important pour vous deux.
            \\
            N'hésitez pas à me faire savoir s'il y a quelque chose que je puisse faire pour vous aider à organiser ou à préparer la réception.

        \subsection*{C1}
            Je vous remerci sincèrement pour votre invitation, et je vous adresse toutes mes félicitations à l'occasion de votre mariage.
            C'est une très heureuse nouvelle, et je serai ravi d'être parmi vous pour célébrer cet événement si important.
            \\
            J'accepte votre invitation avec grand plaisir et je viendrai accompagné de mon ami.
            Nous serons très heureux de partager cette belle journée avec vous.
            \\
            J'aurais toutefois quelques petites questions concernant l'organisation: 
            souhaitez-vous que les invités arrivent un peu avant la cérémonie?
            Y a-t-il une tenue particuulière à prévoir pour le repas au restaurant? 
            Enfin, avez-vous besoin d'un coup de main pour les préparatifs ou pour l'organisation ce jour-là?
            \\
            Encore toutes mes félicitations à vous deux, et au plaisir de vous revoir très bientôt.
            \\
            Bien amicalement,

\chapter{DELF B1}
    \section{Exemple 1}
        \subsection{Essai}
            Dans les vinngt derinères années, mon pays d'origine, la Chine, a changé beaucoup, 
            par exemple, la augmentation dans le nombre de véhicules électriques, l'amélioration de la technologie et 
            les changements environmentaux.
            \\
            Au cours des dix dernières années, la Chine a été concentrer les resources scientifiques et financières sur l'industrie de mobile électrique,
            et il a devenue l'un des payers leaders dans ce secteur. Comme Tesla dans les États-Unis, 
            il existe de nombreuses marques pour véhicules électriques, l'un des plus grands être BYD. 
            Je suis rétourné en Chine l'été 2023, après obtenir son diplôme d'études secondaires. 
            Presque la moitié des véhicules sur la route étaient électriques. Mon cousin a acheté un grand SUV électrique, nommé Haval.
            L'intérieure est très grand, il y a une écran géant devant la voiture, il est capable d'afficher de la vidéo provenant d'un portable via Bluetooth.
            Je dirais que c'est une changement positif, parce que il y a moins de bruit et de émissions de gaz à effet de serre.
            \\
            La Chine avait fait pleusir progrès technologiques, le plus grand et le plus influent était la reconnaissance faciale.
            Cette technologie était utiliser à partir de la surveillance au paiement des courses, par conséquent, les caméras sont installées partout,
            qui forme un réseau de surveillance. Bien que cette technologie facilite la vie quotidienne des gens, 
            mais je pense que c'est une changement négative, depuis il est trop facile de suivre quelqu'un. 
            \\
            Au cours des dernières années, la Chine a été le plus grand exportatrice pour les marchandises et la transformation des matières premiéres.
            La majeure partie de l'énergie industrielle provient de la combustion du charbon et du gaz, 
            qui provoque la pollution de l'air, la émissions de gaz à effet de serre, ainsi représente les changements environmentaux,
            où le ceil est grise tout le temps, la rivière et l'océan sont bruns, ne sont plus bleus.
            \\
            En conclusion, il y avait les changments positives et negatives pour la Chine, mais dans l'ensemble, 
            la technologie est évolue en permanence.

        \subsection*{Fixes}
            \begin{itemize}
                \item Dans les vingt dernières années $\to$ Au cours des vingt dernières années.
                \item a changé beaucoup $\to$ a beaucoup changé. 
                \item l'augmentation.
                \item les changements environnementaux.
                \item la Chine a été concentrer $\to$ la Chine s'est concentrée sur $\ldots$
                \begin{itemize}
                    \item se concentrer (focus itself).
                \end{itemize}
                \item les ressources (double s)
                \item l'industrie de mobile électrique $\to$ l'industrie des véhicules électriques.
                \item Il a devenue $\to$ elle est devenue.
                \item Comme aux États-Unis avec Tesla or Comme Tesla aux États-Unis.
                \item dont l'une des plus grandes est BYD
                \begin{itemize}
                    \item dont means ``of which, among which, including''
                    \item une not un because ``manque'' is feminine
                \end{itemize}
                \item retourné not ré $\dots$
                \item après obtenir son diplôme $\to$ après avoir obtenu mon diplôme.
                \begin{itemize}
                    \item après avoir + past participle (after having done)
                \end{itemize}
                \item L'intérieure est très grand $\to$ L'intérieur était très spacieux.
                \item moins de bruit et de émissions $\to$ moins de bruit et moins d'émissions.
                \item avait fait plusieurs progrès $\to$ avait fait plusieurs progrès, or a connu plusieurs progrès technologiques.
                \item était la reconnaissance faciale $\to$ le plus important était la reconnaissance faciale.
                \begin{itemize}
                    \item progrès is masculine.
                \end{itemize}
                \item Cette technologie était utilisée.
                \item à partir de la surveillance au paiement (awkward) $\to$ de la surveillance au paiement, pour la surveillance comme pour le paiement.
                \item qui forme un réseau $\to$ ce qui crée un réseau de surveillance.
                \item No mais: Bien que cette technologie facilite la vie quotidienne, je pense que...
                \item un changement négatif
                \item depuis il est trop facile $\to$ car il est trop facile, puisqu'il est trop facile.
                \item le plus grand exportatrice $\to$ le plus grand exportateur, l'une des plus grandes exportatrices. Or simplest: La Chine est devenue l'un des plus grands exportateurs de marchandises.
                \item des changements positifs et négatifs.
                \item la technologie est évolue $\to$ la technologie évolue, la technologie a évolué.
            \end{itemize}

        \subsection*{C1}
            Au cours des vingt dernières années, la Chine a connu des transformations profondes.
            À mes yeux, les plus marquantes sont l'essor des véhicules électriques, le développement des technologies de reconnaissance faciale
            et l'aggravation de certains problèmes environnementaux.
            \\
            D'une part, l'augmentation du nombre de véhicules électriques représente une évolution positive. 
            Grâce à d'importants investissements, la Chine est devenue l'un des acteurs majeurs de ce secteur.
            Cette transition permet de réduire le bruit en ville et, dans une certaine measure, les émissions polluantes.
            \\
            D'autre part, la reconnaissance faciale s'est imposée dans de nombreux aspects de la vie quotidienne,
            de la surveillance au paiement. Même si elle facilite certaines tâches, elle soulève aussi de vraies inquiétudes en matière
            de vie privée et de contrôle social.
            \\
            Enfin, l'industrialisation rapide du pays a eu des conséquences négatives
            sur l'environnement, notamment sur la qualité de l'air.
            \\
            En somme, ces changements montrent à la fois les progrès impressionnants de la Chine
            et les défis qu'elle doit encore relever.

    \newpage
    \section{Exemple 2}
        \subsection{Essai}
            Il y a beaucoup des gens qui vie dans les grandes métropoles, la population se compte souvent en millions, voire en millards.
            Ainsi, nettoyage quotidienne ne résoudra pas le problème à l'origin.
            On peut, et on doit, réduire la pollution dans les grandes métropoles sous deux aspects suivant: la planification urbaine et sensibiliser le public à la protection de la environment.
            \\
            La planification urbaine signifie la planification minutieuse de la système d'égouts qui assure l'eaux usées s'écoule doucement et plus important encore élimine la possibilité de l'inondation intérieur.
            Cela gardera la ville propre. Un exemple du passé est Londres à la fin du XIXe siècle, où les rues étaient sales, jonchées d'ordures à l'odeur épouvantable.
            La cause principe était la système d'egouts étroites, donc l'eaux usées ne s'écoulait pas dans le temps.
            \\
            Même si une grande métropole a une la efficace planification urbaine bien construit. 
            Le sensibiliser de citoyenne est l'aspect plus important. 
            Prenons Singapour comme exemple extrême, il y a beaucoup de nombre de poubelles dans les rues,
            des lois comme l'interdiction du chewing-gum. Le citoyen possède une conscience élevée pour la protection de la environment, 
            donc la ville de Singapour peut être proper comme un jardin.
            Le même case s'applique au japonais citoyens.
            \\
            Enfin, pour garder les grandes métropoles propre, c'est un responsabilité des citoyens et des urbanistes. C'est un effort collaborative.

        \subsection*{Fixes}
            \begin{itemize}
                \item la population s'y compte souvent en millions, voire en milliards.
                \item Ainsi, un nettoyage quotidien ne résoudra pas le problème à l'origine.
                \item sous deux aspects suivant $\to$ sous deux aspects suivants, or de deux maières principales.
                \item qui assure l'eaux usées s'écoule $\to$ qui assure que les eaux usées s'écoulent correctement.
                \item élimine la possibilité de l'inondation intérieur $\to$ réduit les risques d'inondations en ville, or élimine le risque d'inondations dans les zones urbaines.
                \item La cause principe $\to$ la cause principale.
                \item la système d'egouts étroites $\to$ le système d'égouts était trop étroit.
                \item Même si une grande métropole dispose d'une planification urbaine efficace, cela ne suffit pas.
                \item La sensibilisation des citoyens.
                \item l'aspect le plus important (superlative is: le plus important)
                \item Beaucoup de nombre de poubelles $\to$ un grand nombre de poubelles.
                \item Le citoyen possède une conscience élevée $\to$:
                \begin{itemize}
                    \item les citoyens ont une forte conscience écologique.
                    \item les citoyens sont très sensibles à la protection de l'environnement.
                \end{itemize}
                \item Peut être proper comme un jardin $\to$ peut être propre, presque comme un jardin.
                \item Le même case s'applique au japonais citoyens $\to$ Le même constat s'applique aux citoyens japonais.
                \item un effort collectif.
                \item une responsabilité.
            \end{itemize}

        \subsection*{C1}
            Dans les grandes métropoles, la réduction de la pollution représente un défi essentiel en raison de la forte densité de population et de l'intensité des activités humaines.
            À mon avis, deux leviers principaux permettent d'agir efficacement: une planification urbaine rigoureuse et une véritable sensibilisation de la population.
            \\
            D'une part, les pouvoirs publics doivent investir dans des infrastructures adaptées, notamment dans des transports en commun performants, des systèmes d'assainissement modernes et davantage d'espaces verts.
            De telles mesures permettent de limiter la circulation automobile, de mieux gérer les déchets et d'améliorer la qualité de vie.
            \\
            D'autre part, la lutte contre la pollution dépend aussi du comportement des citoyens. 
            Sans discipline collective, même les meilleures politiques urbaines restent insuffisantes.
            Singapour en est un bon exemple: grâce à des règles strictes et à une forte conscience civique, la ville est réputée pour sa propreté.
            \\
            En somme, seule une action commune des autorités et des habitants permettra de réduire durablement la pollution urbaine.

\chapter{Writing Practice}

    \section{Day 1-3}
        \subsection{Fait-il interdire le téléphone à l'école?}
            Je pense que le téléphone devrait être interdire à l'école car en n'apporte aucun avantages à les étudiants,
            les raisons sont: D'abord, en ne bonne santé pour l'oeuil des étudiant. Le téléphone serait en hausse la chance de la myopie.
            Le seconde raison est les étudiants jeune ne sont besoin de le téléphone pour assister avec leur cours ou rien. Ils devraient passer le temps avec leur amis,
            renforcer leur compétences sociales, comment interagir avec les autres.
            En fin, le téléphone est un fardeau pour les étudiant, parce qu'ils sont trop jeune pour avoir le téléphone et en utiliser responsable. \\
            Comme nous tous savions, quand nous regardons l'écran de téléphone, le cristallin de l'œil s'ajuste à une position. Depuis l'écran est très proche à notre œils.
            De temps en temps, le cristallin et le muscle va se fatiguer et le sera difficile de s'adapter. Par conséqquent, les étudiants peut developer la myopie 
            qui eux forcer à apporter les lunettes ou les lentilles, et affecter leur vie tous les jours.\\
            Le cours de l'école est facile et surtout se concentre sur les compétences sociales, logiques et reconnaissances.
            Le téléphone ne viens pas entre en jeu. \\
            Finalement, le téléphone est un objet du valuer pour les étudiantes porter. Ils peuvent le tomber quand faire l'activités, par exemple, courir dans le terrain du joue. \\
            En conclusion, téléphone ne porteraient pas les avantages pour les étudiantes, en fait, il retardera leur croissance et leur apprentissage.

            Fixed Version: \\
            Certes, le téléphone peut être utile en case d'urgence. Cependant, à l'école, il est souvent une source de distraction et son utilisation devrait être limitée aux situations nécessaires.
            Ainsi, je pense que le téléphone devrait être interdit à l'école, parce qu'il n'apporte aucun avantage aux élèves, malgré tous les avantages qu'il apporte dans les monde de adulte. \\
            Tout d'abord, il n'est pas bon pour la santé des yeux. L'usage du téléphone peut augmenter le risque de myopie, surtout chez les jeunes.
            Ensuite, la plupart des élèves n'ont pas besoin de leur téléphone pour suivre les cours. À l'école, l'objectif principal est d'apprendre et de développer des conmpétences sociales et logiques. 
            Les élèves devraient donc passer du temps avec leurs amis, renforcer leurs compétences sociales et apprendre à interagir avec les autres.
            De plus, regarder un écran de près pendant longtemps fatigue les yeux: le cristallin et les muscles dovient s'ajuster en permanence. Par conséquent, certains élèves peuvent développer une myopie,
            ce qui les oblige à porter des lunettes ou des lentilles et peut affecter leur quotidienne.
            Enfin, le téléphone peut être un fardeau: certains élèves sont trop jeunes pour l'utiliser de manière responsable.
            Il peut aussi être perdu ou cassé pendant les activités sportives, par exemple lorsqu'on court dans la cour de récréation. \\
            En conclusion, interdire le téléphone à l'école serait utile, car il peut nuire à la santé, distraire les élèves et freiner leur apprentissage. \\

            B2 Level: \\
            Certes, le téléphone peut être utile \textcolor{red}{en cas d'urgence}. 
            Cependant, à l'école, il devient souvent une source de distraction et son utilisation devrait être limitée aux situations vraiment nécessaires.
            Ainsi, je pense que le téléphone devrait être interdit à l'école, car il n'apporte \textcolor{red}{que peu d'avantages aux élèves}, même s'il est très utile dans le monde des adultes. \\
            Tout d'abord, il n'est pas bon pour la santé des yeux. 
            L'usage prolongé du téléphone peut augmenter le risque de myopie, surtout chez les jeunes. 
            En effet, lorsque l'on regarde un écran de près pendant longtemps, le cristallin et les muscles de l'œil doivent s'ajuster en permanence, ce qui fatigue la vue.
            \textcolor{red}{Par conséquent}, certains élèves peuvent développer une myopie, ce qui peut les obliger à porter des lunettes ou des lentilles et affecter leur vie quotidienne. \\
            Ensuite, la plupart des élèves n'ont pas besoin de leur téléphone pour suivre les cours. 
            À l'école, l'objectif principal est d'apprendre et de développer des compétences sociales et logiques. 
            Les élèves devraient plutôt passer du temps avec leurs amis et apprendre à interagir avec les autres. \\
            Enfin, le téléphone peut aussi être un fardeau: certains élèves sont trop jeunes pour l'utiliser de manière responsable,
            et il peut être perdu ou cassé pendant les activités sportives. \\
            En conclusion, interdire le téléphone à l'école aiderait à protéger la santé et à améliorer la concentration.

\end{document}